\section{Choix d'index : l'alternative mesurée en local}\label{sec:mesures}
Nous comparons une base avec le seul index automatique \code{_id} et une base avec des index adaptés aux filtres. Les deux index F1 associent
respectivement zone/date de départ et zone/date d'arrivée. L'index F4 associe
origine, destination et date de départ : les deux égalités précèdent l'intervalle.
Ils ne remplacent pas la lecture des flags et métriques dans les documents.

Nos essais du 5 novembre utilisent janvier complet, 2\,964\,624 lignes,
MongoDB 8.0.16 et Spring Boot 3.5.7 / Temurin 21.0.9 sur un serveur Linux.
Le cgroup est limité à deux CPU et 8 Gio ; API et MongoDB partagent CPU, mémoire
et disque. Pour chaque requête et configuration : un échauffement, puis trois
appels mesurés, sans purge des caches. Les médianes ci-dessous excluent
l'échauffement ; les temps individuels, écarts-types, réponses et plans
\code{explain("executionStats")} sont archivés \cite{mesures}.

\begin{center}\small
\begin{tabular}{llrrr}
\toprule
Service & Requête & Sans index & Index F1 & F1 + F4\\
\midrule
F2 & Zone 181, janvier & 2,313 & 0,209 & ---\\
F2 & Brooklyn, janvier & 3,464 & 1,330 & ---\\
F4 & 132 $\to$ 161 & 2,816 & 0,314 & 0,045\\
F4 & 161 $\to$ 236 & 2,515 & 0,433 & 0,186\\
F5 & Zone 132 & 3,653 & 0,971 & 1,046\\
F5 & Brooklyn & 3,799 & 0,572 & 0,297\\
\bottomrule
\end{tabular}
\end{center}
Les temps sont en secondes. F4 compare les 1--14 et 15--31 janvier ; F5 compare
les semaines du 15 et du 22 janvier avec la baseline du 1er au 22. Les deux
requêtes de F2 et les deux paires F4 comparent la conception sur plusieurs accès
représentatifs ; elles ne remplacent pas E1--E3.

Pour F2 zone 181, les documents examinés passent de 2\,964\,624 à 7\,924.
Pour F4 132 $\to$ 161, ils passent de 111\,942 avec F1 à 2\,459 avec F4.
L'index de paire apporte donc un gain sur cet accès, au coût d'environ
54,1 Mio d'index supplémentaires et 7,46 s de construction dans cet essai.
Le surcoût pendant l'import n'y est pas mesuré. Pour F5, F1 et F1+F4 examinent
les mêmes 133\,084 documents sur la zone 132 : les temps seuls n'établissent pas
un bénéfice causal de F4. Nous retenons les deux index F1 et l'index F4, qui réduisent fortement les documents lus pour un coût de stockage et de construction limité ; nous n'avons retenu aucun index F5 supplémentaire.

Dans ces essais, le contrôle ville de F3 est un scan complet : les médianes
16,30/14,18 s pour la zone 132 sans/avec F1 montrent le gain limité du filtre
local lorsque le reste domine. L'index temporel \code{pickup_datetime_zone}, nécessaire au placement par période, sert aussi ce contrôle. \code{urbanflow.deploy} crée les quatre index avant l'import sur tous les déploiements.
